% 区域、边界分段与外法向。左侧边界段 Gamma_D 规定场值，右侧边界段 Gamma_N 规定法向通量；
% 右侧小图示意 Robin 条件：向外热流与边界温度和环境温度之差成正比。
\begin{tikzpicture}[line width=.75pt,font=\fontsize{8}{10}\selectfont,>=Stealth,
  note/.style={FigureMuted,font=\fontsize{7.5}{9}\selectfont}]
% domain
\fill[FigureBlueFill!45!white] plot[smooth cycle,tension=.8] coordinates
  {(0,0) (1.6,-.55) (3.4,-.25) (4.3,.9) (3.6,2.2) (1.9,2.55) (.3,2.0) (-.45,1.0)};
% Gamma_D: left part, Gamma_N: right part (same smooth curve split in two)
\draw[FigureBlue,line width=1.6pt] plot[smooth,tension=.8] coordinates
  {(1.9,2.55) (.3,2.0) (-.45,1.0) (0,0) (1.6,-.55)};
\draw[FigureGreen,line width=1.6pt] plot[smooth,tension=.8] coordinates
  {(1.6,-.55) (3.4,-.25) (4.3,.9) (3.6,2.2) (1.9,2.55)};
\node[FigureInk] at (1.9,1.1) {$\Omega$};
\node[FigureInk] at (1.9,.6) {$-\Delta u=f$};
\node[FigureBlue,anchor=east] at (-.55,1.6) {$\Gamma_D$：$u=g_D$};
\node[FigureGreen,anchor=west] at (4.35,1.7) {$\Gamma_N$：$\partial_{\bm n}u=g_N$};
% outward normal at a point of Gamma_N
\fill[FigureInk] (4.28,.75) circle (1.2pt);
\draw[->,FigureStroke] (4.28,.75) -- (5.15,.45) node[right,FigureInk] {$\bm n$};
\node[note,anchor=north] at (1.9,-.75) {法向导数 $\partial_{\bm n}u=\nabla u\cdot\bm n$，$\bm n$ 为单位外法向};
% Robin sketch
\begin{scope}[xshift=7.1cm,yshift=.2cm]
  \fill[FigureBlueFill] (0,0) rectangle (1.4,2.0);
  \draw[FigureYellow,line width=1.6pt] (1.4,0) -- (1.4,2.0);
  \node[FigureInk] at (.7,1.0) {$u$};
  \node[FigureInk,anchor=west] at (2.0,1.0) {$u_{\mathrm{env}}$};
  \draw[->,FigureRed,line width=.9pt] (1.15,1.55) -- (2.15,1.55);
  \node[note,anchor=south west] at (1.5,1.62) {热流};
  \node[FigureInk,anchor=north] at (1.4,-.12) {$\kappa\partial_{\bm n}u+\alpha u=\alpha u_{\mathrm{env}}$};
  \node[note,anchor=north] at (1.4,-.55) {Robin：通量与温差成正比};
\end{scope}
\end{tikzpicture}
